% SPDX-License-Identifier: CC-BY-4.0
% Copyright (c) 2026 Martin Schröder <info@swedishembedded.com>

\section{Reproducing the Vim experiment step by step}\label{app:repro-vim}

\begin{sloppypar}
This appendix is the procedure a third party runs, from a fresh machine, to repeat the Vim source-to-transfer experiment
(\cref{sec:introduction}) from the published material. Every command block below is meant to be pasted, in order, into one
\texttt{bash} session; each block's first line names its part: \texttt{[setup]} (build, fetch, prepare), \texttt{[serve]}
(the model servers), \texttt{[smoke]} (a check of the whole chain in minutes), \texttt{[full]} (the experiment) and
\texttt{[analysis]}. The revisions it pins are those of \cref{tab:artifacts,tab:builds}. The procedure was run as written,
\texttt{[setup]}, \texttt{[smoke]} and \texttt{[analysis]} end to end, on 2026-11-27 (\cref{app:repro-vim-test}).
\end{sloppypar}

\subsection{What is published and what each step needs}\label{app:repro-vim-scope}

\begin{center}\footnotesize
\begin{tabular}{>{\raggedright\arraybackslash}p{38mm}>{\raggedright\arraybackslash}p{56mm}>{\raggedright\arraybackslash}p{50mm}}
\toprule
Step & Inputs & Reproducible from the published material \\
\midrule
Tasks from the manual (forward and outcome-first authoring, review, repair) & the manual, the authoring operators, prompts and
schemas & no: the authoring pipeline is not published; its output, the 1026 forward and 467 outcome-first tasks, is, and is
checked byte for byte against the runs' pinned tasks (step~\ref{app:repro-vim-verify}) \\
Oracle admission & published tasks & yes \\
Student and teacher attempts & published tasks, Qwen3-4B, Qwen3.8-27B-FP8 & yes (new samples, not the same trials) \\
Arm datasets B, C, D, F, G & attempts of the step above & yes; the exact datasets the arms trained on are also published \\
Training & an arm dataset & yes, from rebuilt or from the published datasets \\
Exam evaluation (T2/T4) & 100 exams, candidate and base & yes: the exams, their checks and reference answers are public \\
T1 evaluation & held-out authored tasks & yes \\
Registered analysis & per-task outcomes & yes for a new run; for the published run only its aggregates \\
\bottomrule
\end{tabular}
\end{center}

\begin{sloppypar}
The exams are the evaluation side of the public dataset \nolinkurl{swedishembedded/vim-headless-mastery}: each exam's
instruction, starting files, checker (\texttt{tests/test.sh} with the expected files) and reference answer are published, so a
third party grades its own candidates exactly as the experiment did. At revision \texttt{ff83d49b} the published exams carry an
older copy of the environment's three files (the image recipe, the vim tool and its configuration) than the run used; the
procedure lays the run's copies, which every published practice task carries, over them, after which all 1263 exam files are
byte-identical to the ones the run pinned (step~\ref{app:repro-vim-verify}). What a third party cannot recompute is the
published run's own per-exam outcomes: they are item-level exam data kept with the examiner's private evidence, and the
published run is checked through its aggregates (the evidence dataset's \texttt{evaluation/} and \texttt{decision/} files,
\cref{tab:repro-vim-outputs}).
\end{sloppypar}

\subsection{Prerequisites and hardware}\label{app:repro-vim-prereq}

\begin{sloppypar}
\textbf{Software.} Linux (the run: Ubuntu 26.04 on aarch64, in a container on a cluster node) with unprivileged user
namespaces or a Docker daemon (the sandbox of each attempt; \texttt{unshare -Ur true} must succeed); git with git-lfs; a Rust
toolchain (the run: rustc 1.100.0; \texttt{rustup} stable); Python~3.12 or later with \texttt{venv} (the run: 3.14);
\texttt{zstd}; the \texttt{hf} command-line client, logged in (\texttt{hf auth login}) with read access to the private
\texttt{swedishembedded} datasets; for the model servers the CUDA toolkit~13 (\texttt{nvcc} on \texttt{PATH},
\texttt{CUDA\_PATH} set) and an NVIDIA driver; network access to github.com, huggingface.co, the Debian snapshot archive and
the Docker registry (each task image is built from its pinned \texttt{debian:trixie-slim} digest and a fixed snapshot of
the archive). Vim is not needed on the host: every Vim the experiment runs is the task image's (Debian's 9.1.1230). Sources:
brain is public (\nolinkurl{github.com/swedishembedded/brain}); splinter and sven at the commits used are in the authors'
repositories \nolinkurl{github.com/mkschreder/splinter} and \nolinkurl{github.com/mkschreder/sven}, available on request.
\end{sloppypar}

\begin{sloppypar}
\textbf{Hardware.} The run used four nodes of four NVIDIA GH200 (96~GB of device memory each, 288 Arm cores per node): one
brain server per card (16), each node's attempts and evaluations through a router over its cards, each arm-seed's training on
one card. The weights take 7.5~GiB (Qwen3-4B, bf16) and 28.8~GiB (Qwen3.8-27B-FP8); a machine with fewer or smaller cards
was not tried. Disk: see \cref{app:repro-vim-test} for what the test used.
\end{sloppypar}

\subsection{Sources and builds}\label{app:repro-vim-build}

{\scriptsize
\begin{verbatim}
# [setup] Where everything goes, the pinned revisions, and where the model servers are.
export REPRO=${REPRO:-$HOME/repro}
mkdir -p "$REPRO/tmp" "$REPRO/cache" && chmod 700 "$REPRO" && cd "$REPRO"
export TMPDIR=$REPRO/tmp CARGO_BUILD_JOBS=${CARGO_BUILD_JOBS:-48}
X=$REPRO/experiment   # the task repository the recipes run in
SPLINTER_GIT=${SPLINTER_GIT:-https://github.com/mkschreder/splinter.git}
SPLINTER_REV=bde8e09f94bb584796acfd9e792db788ba54d263
SVEN_GIT=${SVEN_GIT:-https://github.com/mkschreder/sven.git}
SVEN_REV=01eabb507058cb37d0254a0340cb4dfe87073576
BRAIN_GIT=https://github.com/swedishembedded/brain.git
BRAIN_REV=d6bd3baaff537ed69b9fdb23b2f0c702b43f7658
# The model servers of [serve]; a machine that has servers elsewhere sets these first.
export BRAIN_KEY_FILE=${BRAIN_KEY_FILE:-$REPRO/brain.key}
export BRAIN_CHAT_URL=${BRAIN_CHAT_URL:-http://127.0.0.1:8990/v1}
export BRAIN_TRAIN_URL=${BRAIN_TRAIN_URL:-http://127.0.0.1:8813/v1}
export BRAIN_CA_FILE=${BRAIN_CA_FILE:-}
export BRAIN_ARTIFACTS=${BRAIN_ARTIFACTS:-$REPRO/brain-store/artifacts}
export BRAIN_ADAPTERS=${BRAIN_ADAPTERS:-$REPRO/brain-store/adapters}
# What the builds below make, on PATH; the worker starts the agent from SPLINTER_SVEN_BIN.
export PATH=$REPRO/splinter/target/release:$REPRO/splinter/workers/harbor/.venv/bin:$PATH
export PATH=$REPRO/brain/target/release:$PATH SPLINTER_SVEN_BIN=$REPRO/sven/target/release/sven
\end{verbatim}
}

\begin{sloppypar}
Brain \texttt{d6bd3baa} is the published history's copy of \texttt{b726c045b}, the build every server of the run ran (a later
rewrite of the history inserted a version bump; the two differ only in the version number, 7.0.0 for 6.0.0, and the lock
file's entries for it). Sven \texttt{01eabb50} is the agent every attempt ran (\texttt{sven 5.9.2}); splinter's own lock pins
the sven and brain library code it links. The run's controllers were splinter \texttt{cc5581333} (both material runs) and
\texttt{f6661e34} (the arms); \texttt{SPLINTER\_REV} is the later commit this procedure was tested with, which holds every
option the arms used.
\end{sloppypar}

{\scriptsize
\begin{verbatim}
# [setup] Clone at the pinned commits and build.
git clone -q "$BRAIN_GIT" brain && git -C brain -c advice.detachedHead=false checkout -q $BRAIN_REV
git clone -q "$SVEN_GIT" sven && git -C sven -c advice.detachedHead=false checkout -q $SVEN_REV
git clone -q "$SPLINTER_GIT" splinter
git -C splinter -c advice.detachedHead=false checkout -q $SPLINTER_REV
make -C sven build/release JOBS=$CARGO_BUILD_JOBS
make -C splinter build
make -C splinter worker/install
\end{verbatim}
}

{\scriptsize
\begin{verbatim}
# [serve-build] Brain, on the machine with the GPUs (CUDA 13 on PATH, CUDA_PATH set).
make -C "$REPRO/brain" build/release
\end{verbatim}
}

\subsection{Data at pinned revisions}\label{app:repro-vim-data}

\begin{sloppypar}
The task datasets hold tens of thousands of small files; a git clone fetches each in seconds, where per-file downloads run into
the Hub's request limit. The credential helper hands git the \texttt{hf} login's token without writing it anywhere.
\end{sloppypar}

{\scriptsize
\begin{verbatim}
# [setup] Fetch the five repositories at the revisions the paper reports.
hfgit() { git -c credential.helper= -c 'credential.helper=!f() { echo username=hf;
  echo "password=$(hf auth token 2>/dev/null)"; }; f' \
  -c filter.lfs.process='git-lfs filter-process' -c filter.lfs.required=true "$@"; }
hfclone() { hfgit clone -q --no-checkout "https://huggingface.co/$1" "$3" &&
  hfgit -C "$3" -c advice.detachedHead=false checkout -q "$2"; }
mkdir -p data && cd data
hfclone datasets/swedishembedded/vim-headless-mastery \
  ff83d49b1b7d8636300eef73bcc61377af8829b1 vim-headless-mastery
hfclone datasets/swedishembedded/vim-headless-mastery-manual \
  3636839122fe6369d023e8d959576f30074a0373 vim-headless-mastery-manual
hfclone datasets/swedishembedded/vim-headless-mastery-sft-manual \
  bc55f91f923fc5bdc02b398e1eeebb75da421360 vim-headless-mastery-sft-manual
hfclone datasets/swedishembedded/source-to-transfer-evidence \
  413259d233569d0144154e1e2e5198e99095fdc4 source-to-transfer-evidence
hfclone swedishembedded/vim-manual-transfer-adapters \
  96e33e51d1e456b104e9c142d11f03ceb8e4ca3d vim-manual-transfer-adapters
cd "$REPRO"
\end{verbatim}
}

\begin{sloppypar}
\nolinkurl{vim-headless-mastery-manual} holds the practice tasks (\nolinkurl{train-forward/}, 1026; \nolinkurl{train-backward/}, 467);
\nolinkurl{vim-headless-mastery} the course, whose \nolinkurl{eval/} side is the 100 exams; \nolinkurl{source-to-transfer-evidence} the
arm datasets exactly as trained, the runs' recipes, locks and revisions, every material stage's product manifest, the aggregates
and the analysis code; \nolinkurl{vim-manual-transfer-adapters} the 15 adapters with their training reports;
\nolinkurl{vim-headless-mastery-sft-manual} the same trajectories as chat records for other trainers (arm B's and G's seed-1 train
and validation records with repeated solutions moved to \nolinkurl{duplicates}, the T1 records in \nolinkurl{heldout}, and
policy-gradient records). The procedure below does not need the last one.
\end{sloppypar}

\subsection{The task repository}\label{app:repro-vim-repo}

\begin{sloppypar}
Splinter runs a recipe inside a task repository: the tasks, the collections that name them, the sandbox's environment, the
evaluation plans and the files the agent is given. The task identifiers must be the run's (\texttt{train/manual-\ldots},
\texttt{train/backward-\ldots}, \texttt{eval/exam-\ldots}), since the T1 hold-out is a hash of the seed and the task identifier;
the layout below gives them.
\end{sloppypar}

{\scriptsize
\begin{verbatim}
# [setup] The tasks, with the run's environment files laid over the published exams.
mkdir -p $X/tasks/forward $X/tasks/backward $X/tasks/course $X/collections $X/eval \
  $X/environments/vim $X/presets $X/prompts $X/recipes $X/datasets
cp -r data/vim-headless-mastery-manual/train-forward $X/tasks/forward/train
cp -r data/vim-headless-mastery-manual/train-backward $X/tasks/backward/train
cp -r data/vim-headless-mastery/eval $X/tasks/course/eval
env_src=$X/tasks/forward/train/manual-0007c64a-update-menu-docs/environment
for exam in $X/tasks/course/eval/*/; do
  cp $env_src/Dockerfile $env_src/vim_mcp.py $env_src/sven.yaml "$exam/environment/"
done
for name in forward backward course; do
  printf 'schema: splinter.collection/v1\nname: %s\nroot: tasks/%s\n' $name $name \
    > $X/collections/$name.yaml
done
\end{verbatim}
}

\paragraph{Check against the run.}\label{app:repro-vim-verify}
The material runs' manifests in the evidence dataset list every file of the tasks they pinned with its SHA-256; this compares
the tree just built with them.

{\scriptsize
\begin{verbatim}
# [setup] Every task file byte-identical to what the runs pinned (prints three lines).
mkdir -p $REPRO/evidence && tar -C $REPRO/evidence --zstd \
  -xf data/source-to-transfer-evidence/material/material-manifests.tar.zst
cat > $REPRO/verify-tasks.py <<'EOF'
import hashlib, json, os, sys
evidence, tasks = sys.argv[1], sys.argv[2]
for run, stage, tree in [("material2", "tasks", "forward"), ("backward2", "tasks", "backward"),
                         ("material2", "exams", "course")]:
    pinned = json.load(open(f"{evidence}/{run}/{stage}/manifest.json"))["files"]
    want = {f["path"]: f["digest"] for f in pinned if "/" in f["path"]}
    have = {}
    for d, _, names in os.walk(os.path.join(tasks, tree)):
        for n in names:
            p = os.path.join(d, n)
            have[os.path.relpath(p, os.path.join(tasks, tree))] = \
                "sha256:" + hashlib.sha256(open(p, "rb").read()).hexdigest()
    same = sum(have.get(k) == v for k, v in want.items())
    gone = sorted({k.split("/")[1] for k in want if k not in have})
    print(f"{tree}: {same} of {len(want)} pinned files identical, {len(have) - same} other files;"
          f" tasks not published: {gone}")
    if len(have) != same or any(k in have and have[k] != v for k, v in want.items()):
        sys.exit(f"{tree}: the tree differs from the run's")
EOF
python3 -I $REPRO/verify-tasks.py $REPRO/evidence $X/tasks
\end{verbatim}
}

\begin{sloppypar}
Expected: \texttt{forward: 12709 of 12757 pinned files identical, 0 other files}, naming the four forward tasks the oracle
did not solve and that are not published (their expected files were recorded wrongly, \cref{sec:r-admission});
\texttt{backward: 5890 of 5890}; \texttt{course: 1263 of 1263}.
\end{sloppypar}

{\scriptsize
\begin{verbatim}
# [setup] The sandbox, the plans, the agent's files and the machine's connection.
cat > $X/environments/vim/environment.yaml <<'EOF'
schema: splinter.environment/v1
isolation: confined
sandbox:
  resources: { cpus: 2, memory_gb: 4, timeout: 10m }
network:
  mode: model-gateway-only
EOF
for side in exam:eval unseen:train; do
  cat > $X/eval/vim-${side%%:*}-transfer.yaml <<EOF
schema: splinter.evalplan/v1
split: ${side##*:}
environment: environments/vim/environment.yaml
attempts: 4
concurrency: 64
seed: 20261126
decision: { alpha: 0.05, min_difference: 0.0 }
EOF
  sed 's/^attempts: 4$/attempts: 1/; s/^concurrency: 64$/concurrency: 8/' \
    $X/eval/vim-${side%%:*}-transfer.yaml > $X/eval/smoke-${side%%:*}.yaml
done
cat > $X/presets/sven-no-thinking.yaml <<'EOF'
providers:
  gateway:
    driver_options:
      chat_template_kwargs:
        enable_thinking: false
model:
  sampling_best_effort: true
EOF
echo 'Before each tool call, state in one or two short sentences what you observed and what you' \
  'will do next. No long deliberation.' > $X/prompts/narrate-before-each-call.md
printf 'connections:\n  brain: { protocol: openai, url: "http://127.0.0.1:8990/v1" }\n' \
  > $X/connections.yaml
printf 'connections:\n  brain: { protocol: openai, url: "%s", token: "env:SPLINTER_BRAIN_KEY" }\n' \
  "$BRAIN_CHAT_URL" > $REPRO/machine.yaml
\end{verbatim}
}

\begin{sloppypar}
The system prompt's file is byte-identical to the run's: the narration is appended to sven's own prompt in every attempt
and every evaluation, and the training records carry it. The sven preset makes the student answer without a reasoning block
and send its sampling parameters; the teacher ran without it, in its default thinking mode, its reasoning removed from the
records.
\end{sloppypar}

\subsection{The recipes}\label{app:repro-vim-recipes}

\begin{sloppypar}
One generator writes every recipe, so the stages a fork shares with its parent are the same text and are reused, never
recomputed. A recipe is the material of one task family (the pinned tasks, the oracle's admission, eight student attempts and
four teacher attempts per admitted task, and the T1 split: 10\% of the tasks by a hash of seed 41 and the task's identifier),
optionally followed by the datasets of every arm of that family (\texttt{datasets}) or by one arm-seed: its dataset, a LoRA
trained from the base and its evaluations on the exams and on T1 against the base in the same stage.
\texttt{SMOKE=1} shrinks every count and binds the teacher role to the student's model; \texttt{DATA=published} makes an
arm train on its published dataset instead of one rebuilt from the attempts (\cref{app:repro-vim-full}).
\end{sloppypar}

{\scriptsize
\begin{verbatim}
# [setup] gen-recipe.sh FAMILY PART [SEED]: FAMILY forward|backward; PART material,
# datasets, or an arm letter (b c d f on forward, g on backward) and its seed.
cat > $X/gen-recipe.sh <<'GEN'
#!/usr/bin/env bash
set -euo pipefail
family=$1 part=$2 seed=${3:-1}
if [ "${SMOKE:-0}" = 1 ]; then
  coll=smoke-$family exams=smoke-course author=brain:Qwen/Qwen3-4B
  s_att=2 t_att=1 steps=4 every=2 plan=smoke
else
  coll=$family exams=course author=brain:Qwen/Qwen3.8-27B-FP8
  s_att=8 t_att=4 steps=800 every=40 plan=vim
fi
student_args='[--config, presets/sven-no-thinking.yaml, --temperature, "0.7",
            --top-p, "0.8", --top-k, "20"]'
agent="agent: sven
    system_prompt: prompts/narrate-before-each-call.md
    budget: { completion_tokens: 16384, turns: 60 }"
attempts() { # ID MODEL ATTEMPTS CONCURRENCY SEED [ARGS]
  echo "- id: $1
  uses: attempt.collect
  inputs: { tasks: \$stages.tasks.tasks, solved_by: [\$stages.oracle.attempts] }
  with:
    $agent
    model: $2
    environment: environments/vim/environment.yaml
    attempts: $3
    concurrency: $4
    seed: $5"
  [ -z "${6:-}" ] || echo "    args: $6
    files: [presets/sven-no-thinking.yaml]"
}
dataset() { # ID OPTIONS
  echo "- id: $1
  uses: dataset.project
  inputs:
    attempts: [\$stages.student-attempts.attempts, \$stages.teacher-attempts.attempts]
    tasks: \$stages.practice.tasks
  with:
    projection: native-sft
    target: student
    max_block: 16384
    validation: { group_by: task, percent: 10 }
    seed: 4$2"
}
options() { case $1 in
  c) printf '\n    per_task: 1\n    prefer: fewest-actions' ;;
  d) printf '\n    supervise: productive' ;;
  f) printf '\n    control: shuffled-requests' ;;
esac; }
course=""; [ "$part" = material ] || [ "$part" = datasets ] ||
  course="
  course: { kind: files, path: collections/$exams.yaml }"
[ "${DATA:-attempts}" = attempts ] || course="$course
  records: { kind: files, path: datasets/arm-$part }"
cat <<EOF
schema: splinter.recipe/v2
name: vim-manual-$family-$part-s$seed
inputs:
  tasks: { kind: files, path: collections/$coll.yaml }$course
models:
  author: $author
  student: brain:Qwen/Qwen3-4B
stages:
- id: tasks
  uses: tasks.pin
  inputs: { collection: \$inputs.tasks }
  with: { split: train }
- id: oracle
  uses: attempt.collect
  inputs: { tasks: \$stages.tasks.tasks }
  with: { agent: oracle, environment: environments/vim/environment.yaml, attempts: 1,
          seed: 1, concurrency: 64 }
$(attempts student-attempts student $s_att 512 21 "$student_args")
$(attempts teacher-attempts author $t_att 384 22)
- id: practice
  uses: tasks.pin
  inputs: { collection: \$inputs.tasks }
  with: { split: train, holdout: { percent: 10, seed: 41, side: kept } }
- id: unseen
  uses: tasks.pin
  inputs: { collection: \$inputs.tasks }
  with: { split: train, holdout: { percent: 10, seed: 41, side: held } }
EOF
case $part in
  material) ;;
  datasets) if [ $family = forward ]; then arms="b c d f"; else arms=g; fi
    for a in $arms; do dataset sft-$a "$(options $a)"; done ;;
  *) if [ "${DATA:-attempts}" = attempts ]; then dataset sft "$(options $part)"
    else printf -- '- id: sft\n  uses: dataset.import\n  inputs: { records: $inputs.records }\n'
      printf '  with: { target: student, max_block: 16384 }\n'; fi
    cat <<EOF
- id: lora
  uses: model.train
  inputs: { dataset: \$stages.sft.dataset }
  with:
    base: student
    method: lora
    seed: $seed
    params: { rank: 16, alpha: 32, steps: $steps, grad_accum: 4, learning_rate: 0.0002,
              weight_decay: 0, max_block: 16384, base_precision: bf16,
              eval_every: $every, patience: 4 }
- id: exams
  uses: tasks.pin
  inputs: { collection: \$inputs.course }
  with: { split: eval }
EOF
    for ev in exam:exams unseen:unseen; do cat <<EOF
- id: eval-${ev%%:*}
  uses: evaluation
  inputs: { candidate: \$stages.lora.model, tasks: \$stages.${ev##*:}.tasks }
  with:
    plan: eval/$plan-${ev%%:*}$( [ $plan = vim ] && echo -transfer ).yaml
    $agent
    baseline: student
    args: $student_args
    files: [presets/sven-no-thinking.yaml]
    comparisons: 1
EOF
    done ;;
esac
GEN
chmod +x $X/gen-recipe.sh
\end{verbatim}
}

{\scriptsize
\begin{verbatim}
# [setup] factory NAME RECIPE [PARENT]: run a recipe (or fork PARENT's run) into $RUNS/NAME.
factory() {
  local runs=${RUNS:-$REPRO/runs} verb=(run "$2")
  [ -z "${3:-}" ] || verb=(fork "$runs/$3" --recipe "$2")
  ( cd "$X" && export XDG_CACHE_HOME=$REPRO/cache SPLINTER_BRAIN_KEY=$(cat "$BRAIN_KEY_FILE")
    splinter "${verb[@]}" \
      --out "$runs/$1" --state "$runs/.state/$1" --config "$REPRO/machine.yaml" \
      --brain-endpoint "$BRAIN_TRAIN_URL" ${BRAIN_CA_FILE:+--brain-ca-file "$BRAIN_CA_FILE"} \
      --brain-key-file "$BRAIN_KEY_FILE" --brain-artifact-store "$BRAIN_ARTIFACTS" \
      --brain-adapter-dir "$BRAIN_ADAPTERS" --cluster off --job-launcher local )
}
mkdir -p $REPRO/runs
for f in forward backward; do
  $X/gen-recipe.sh $f material > $X/recipes/$f.yaml
  $X/gen-recipe.sh $f datasets > $X/recipes/$f-datasets.yaml
done
for a in b c d f g; do for s in 1 2 3; do
  f=forward; [ $a = g ] && f=backward
  $X/gen-recipe.sh $f $a $s > $X/recipes/arm-$a-s$s.yaml
done; done
(cd $X && for r in recipes/*.yaml; do splinter validate "$r" --config "$REPRO/machine.yaml" \
  >/dev/null || exit 1; done && echo "every recipe valid")
\end{verbatim}
}

\subsection{Model servers}\label{app:repro-vim-serve}

\begin{sloppypar}
On a fresh machine, the models come from the Hub at the revisions the run's store held, and brain serves one server per card
(here four, ports 8810--8813) behind a router on port 8990 that the attempts and evaluations use; training jobs go to one card's
server directly (a router does not route jobs). With one card, serve one server on port 8990 and set
\texttt{BRAIN\_TRAIN\_URL} to it. Servers run without speculative decoding (\cref{app:deviations}).
\end{sloppypar}

{\scriptsize
\begin{verbatim}
# [serve] Models, a key, four servers and a router (one card: start only gpu0, on 8990).
S=$REPRO/brain-data && mkdir -p $S $BRAIN_ARTIFACTS $BRAIN_ADAPTERS
brain pull --brain-data-dir $S \
  https://huggingface.co/Qwen/Qwen3-4B/tree/1cfa9a7208912126459214e8b04321603b3df60c
brain pull --brain-data-dir $S \
  https://huggingface.co/Qwen/Qwen3.8-27B-FP8/tree/c3a462cc3d38cf351f3b26cbd763469eeb74ef81
[ -s "$BRAIN_KEY_FILE" ] || (umask 077 && head -c 32 /dev/urandom | od -An -tx1 | tr -d ' \n' \
  > "$BRAIN_KEY_FILE")
for n in 0 1 2 3; do
  mkdir -p $S/rt-$n && chmod 700 $S/rt-$n && rm -f $S/ready-$n
  setsid brain serve --openai 127.0.0.1:$((8810 + n)) --device gpu$n --backend cuda --max-batch 64 \
    --speculative off --api-key-file "$BRAIN_KEY_FILE" --training \
    --artifact-store $BRAIN_ARTIFACTS --serve-adapters $BRAIN_ADAPTERS --adapter-targets all \
    --models-dir $S/models --ready-file $S/ready-$n --runtime-dir $S/rt-$n \
    --cache-dir $S/cache-$n > $S/serve-$n.log 2>&1 < /dev/null &
  echo http://127.0.0.1:$((8810 + n))
done > $S/upstreams
mkdir -p $S/rt-router && chmod 700 $S/rt-router
CUDA_VISIBLE_DEVICES= setsid brain serve --openai 127.0.0.1:8990 --upstreams-file $S/upstreams \
  --upstream-key-file "$BRAIN_KEY_FILE" --api-key-file "$BRAIN_KEY_FILE" --models-dir $S/models \
  --runtime-dir $S/rt-router --ready-file $S/ready-router > $S/router.log 2>&1 < /dev/null &
until [ -e $S/ready-0 ] && [ -e $S/ready-1 ] && [ -e $S/ready-2 ] && [ -e $S/ready-3 ] &&
  [ -e $S/ready-router ]; do sleep 5; done
\end{verbatim}
}

\subsection{Smoke check}\label{app:repro-vim-smoke}

\begin{sloppypar}
Before the full run, a few tasks go through every stage: six forward practice tasks and two of the T1 hold-out, three
outcome-first tasks and one of its hold-out, four exams; two student attempts and one ``teacher'' attempt (the student's model)
per task, the datasets of all five arms, and arm B trained for four steps and evaluated once per exam and T1 task, against the
base; then arm B again, four steps on its published dataset.
\end{sloppypar}

{\scriptsize
\begin{verbatim}
# [smoke] A handful of tasks through every stage.
smoke_coll() { # NAME ROOT ID...
  local name=$1 root=$2; shift 2
  { printf 'schema: splinter.collection/v1\nname: %s\nroot: %s\ninclude:\n' $name $root
    printf '  - %s\n' "$@"; } > $X/collections/$name.yaml
}
smoke_coll smoke-forward tasks/forward train/manual-044db1d4-clean-the-log-header \
  train/manual-0c9875ff-remove-trailing-spaces \
  train/manual-10df275c-strip-trailing-spaces-from-log \
  train/manual-18ebaac5-replace-tabs-with-spaces \
  train/manual-2d97698f-remove-first-line-of-notes-md train/manual-2fe2f268-update-navigation-help \
  train/manual-023bed0c-fix-log-timestamps train/manual-11a434a1-fix-config-indentation
smoke_coll smoke-backward tasks/backward train/backward-3a1b4802-reformat-netrw-help-tabs \
  train/backward-b05f6cc7-sort-log-lines-by-timestamp \
  train/backward-f827e63d-remove-trailing-whitespace-in-config-fil \
  train/backward-0570782d-clean-up-vim-plugin-configs
smoke_coll smoke-course tasks/course eval/exam-001-shopping-list eval/exam-002-sensor-readings \
  eval/exam-003-release-checklist eval/exam-004-request-builder
for f in forward backward; do SMOKE=1 $X/gen-recipe.sh $f datasets > $X/recipes/smoke-$f.yaml; done
SMOKE=1 $X/gen-recipe.sh forward b 1 > $X/recipes/smoke-arm-b.yaml
tar -C $X/datasets --zstd -xf data/source-to-transfer-evidence/arm-datasets/arm-b.tar.zst \
  arm-b/train.jsonl arm-b/validation.jsonl
SMOKE=1 DATA=published $X/gen-recipe.sh forward b 1 > $X/recipes/smoke-published-b.yaml
mkdir -p $REPRO/smoke-runs
( export RUNS=$REPRO/smoke-runs   # the smoke's runs, apart from the experiment's
  factory forward recipes/smoke-forward.yaml
  factory backward recipes/smoke-backward.yaml
  factory arm-b-s1 recipes/smoke-arm-b.yaml forward
  factory published-b-s1 recipes/smoke-published-b.yaml forward )
for r in arm-b-s1 published-b-s1; do grep -A4 '^| arm ' $REPRO/smoke-runs/$r/report.md; done
\end{verbatim}
}

\begin{sloppypar}
What the test printed (counts are samples and vary): the oracle solved all 8 forward and 4 outcome-first tasks; the student
passed 12 of 16 and 6 of 8 attempts, the 4B as teacher 6 of 8 and 3 of 4; the datasets of B, D and F held 18 records (15
train, 3 validation) of 6 tasks, C 6 records, G 9; the two T1 tasks were left out of every dataset (\texttt{not\_in\_tasks}),
as in the run. Arm B's four steps took its validation loss from 0.93 to 0.39; on its published dataset the base's validation
loss was 2.4885, the published arm~B seed~1's own (\texttt{lora/report.json}, the same 380 validation records), and 1.06
after four steps. No candidate and no base passed one of the four exams or two T1 tasks.
A run stopped at any point is continued by running the same commands again: every committed stage is reused and only lost
attempts are run again.
\end{sloppypar}

\subsection{The full experiment}\label{app:repro-vim-full}

{\scriptsize
\begin{verbatim}
# [full] Material of both families, every arm's dataset, then 15 arm-seeds as forks.
factory forward recipes/forward.yaml
factory backward recipes/backward.yaml
factory forward-datasets recipes/forward-datasets.yaml forward
factory backward-datasets recipes/backward-datasets.yaml backward
for a in b c d f g; do for s in 1 2 3; do
  parent=forward; [ $a = g ] && parent=backward
  factory arm-$a-s$s recipes/arm-$a-s$s.yaml $parent
done; done
\end{verbatim}
}

\begin{sloppypar}
The arm-seeds are independent: the run started each on its own card (\texttt{BRAIN\_TRAIN\_URL} that card's server), all 15 at
once over 16 cards. Each run directory's \texttt{report.md} (\texttt{splinter report}) gives every stage's counts;
\texttt{splinter inspect RUN --stage STAGE} its products. \Cref{tab:repro-vim-outputs} lists what to compare with the published
run.
\end{sloppypar}

\begin{table}[t]\centering\footnotesize
\caption{What the full procedure produces, what the published run measured (in the evidence file named), and the run's wall
clock on 16 GH200, read from the commit times of its stages' products.}
\label{tab:repro-vim-outputs}
\begin{tabular}{>{\raggedright\arraybackslash}p{30mm}>{\raggedright\arraybackslash}p{80mm}>{\raggedright\arraybackslash}p{34mm}}
\toprule
Product & Published run & Wall clock (run) \\
\midrule
forward material & oracle 1026 of 1026 solved; student 8204 counted attempts, pass rate 0.100, 231 tasks solved; teacher 4095,
0.811, 944 (\texttt{material/funnel.json}) & attempts of both models 5.0~h (14:50--19:48), a restart included \\
outcome-first material & oracle 467 of 467; student 3688, 0.031, 36 solved; teacher 1868, 0.722, 410 & student 33~min,
teacher 32~min \\
arm datasets & B and D 3359 train and 380 validation records of 857 tasks, C 771 train, F as B with the requests deranged, G
1100 train records of 348 tasks (\texttt{arm-datasets/*/composition.json}) & under a minute \\
training & 800 steps, the checkpoint of lowest validation loss of every 40 (arm B seed 1: 2.49 to 0.60;
\texttt{evaluation/aggregates.json}) & 16--51~min per arm-seed on one card \\
exam evaluation & passes of 1200 attempts per arm: B 22, C 10, D 22, F 0, G 42; the paired base 0
(\texttt{decision/decision-eval-exam.json}) & 49--71~min per arm-seed (800 attempts, four arm-seeds per node) \\
T1 evaluation & B 705 of 1128, C 642, D 723, F 9, G 127 of 612 (\texttt{decision/decision-eval-unseen.json}) & not recorded
in the evidence \\
\bottomrule
\end{tabular}
\end{table}

\paragraph{Training on the exact published datasets.}
The attempts are new samples, so rebuilt datasets differ from the published ones in their trials. To train on exactly what the
arms trained on, an arm's recipe imports its published dataset (\texttt{dataset.import}, the records carried over byte for byte) in place
of its \texttt{dataset.project} stage; the three seeds of an arm trained on byte-identical files:

{\scriptsize
\begin{verbatim}
# [full] Arm ARM, seed SEED trained on its published dataset (here arm B, seed 1).
ARM=b SEED=1 F=forward; [ $ARM = g ] && F=backward
tar -C $X/datasets --zstd -xf data/source-to-transfer-evidence/arm-datasets/arm-$ARM.tar.zst \
  arm-$ARM/train.jsonl arm-$ARM/validation.jsonl
DATA=published $X/gen-recipe.sh $F $ARM $SEED > $X/recipes/published-$ARM-s$SEED.yaml
factory published-$ARM-s$SEED recipes/published-$ARM-s$SEED.yaml $F
\end{verbatim}
}

\subsection{Registered analysis, tables and figures}\label{app:repro-vim-analysis}

\begin{sloppypar}
The analysis code is the evidence dataset's. \texttt{export-results.py} reads the per-task outcomes of every arm run's two
evaluations (it names the run's own directory, which the first line below replaces); \texttt{decide.py} makes the registered
decision (per exam, the passes of an arm's 12 attempts over three seeds against its paired base's; one-sided exact sign test,
Holm over B, C, D, G against the base; each arm against F) and \texttt{secondary.py} the secondary comparisons with
task-bootstrap intervals (\texttt{random.Random(20261126)}, 10000 resamples).
\end{sloppypar}

{\scriptsize
\begin{verbatim}
# [analysis] The decision and secondary comparisons of this repetition, and its aggregates.
A=$REPRO/analysis && mkdir -p $A/results && cp data/source-to-transfer-evidence/analysis/*.py $A/
sed -i "s|glob.glob(\"[^\"]*/runs/arm-\*\")|glob.glob(\"${RUNS:-$REPRO/runs}/arm-*\")|" \
  $A/export-results.py
python3 -I $A/export-results.py
for ev in eval-exam eval-unseen; do
  python3 -I $A/decide.py $A/results $ev > $A/decision-$ev.json
  python3 -I $A/secondary.py $A/results $ev > $A/secondary-$ev.json
done
\end{verbatim}
}

\begin{sloppypar}
The paper's tables and figures are generated from the evidence dataset's aggregate files, which \texttt{data/vim-transfer/} in
the paper's directory holds as byte copies. This takes the paper from the current main branch, fetches the evidence revision
its \texttt{data/artifacts.csv} names (later revisions add arms and files to the earlier ones), copies the aggregates over the paper's
and builds the PDF (pdflatex and latexmk). To typeset a repetition's own results, copy its decision and secondary files over
\texttt{paper-evidence/decision/} before \texttt{make vim-data}.
\end{sloppypar}

{\scriptsize
\begin{verbatim}
# [analysis] Tables and figures from the evidence revision the paper pins, and the paper.
git -C splinter fetch -q origin main
git -C splinter worktree add -q --detach $REPRO/paper FETCH_HEAD
P=$REPRO/paper/papers/generalized-continuous-learning
rev=$(awk -F, '$2 == "swedishembedded/source-to-transfer-evidence" { print $4 }' $P/data/artifacts.csv)
hfclone datasets/swedishembedded/source-to-transfer-evidence $rev $REPRO/paper-evidence
make -C $P vim-data EVIDENCE=$REPRO/paper-evidence
make -C $REPRO/paper papers/generalized-continuous-learning/pdf
git -C $REPRO/paper status --short papers/generalized-continuous-learning/data
\end{verbatim}
}

\begin{sloppypar}
An empty status means the paper's data equal the published evidence byte for byte; the PDF is
\texttt{paper/target/papers/generalized-continuous-learning.pdf}. The experiment's code is \texttt{SPLINTER\_REV}'s; the
paper's text, which includes this appendix, is main's.
\end{sloppypar}

\subsection{Not reproducible from the published material}\label{app:repro-vim-limits}

\begin{sloppypar}
Authoring the tasks from the manual (both pipelines: proposing, Vim's admission of each reference, repair, review, the
outcome-first reproduction and request writing) needs the authoring operators, their prompts and output schemas and the
manual's section headings, which are not published; the procedure starts from their published output, which
step~\ref{app:repro-vim-verify} shows identical to the tasks the run used. The published run's per-exam outcomes are not
published (exam item data); its registered decision is checked through \texttt{decision/*.json} and
\texttt{evaluation/aggregates.json}, and a repetition makes its own. The published adapters themselves are not evaluated by a command
here: a repetition trains again on the published datasets (above), and its adapters differ from the published ones in their
weights (in the test, two trainings of the same records on the same server with the same seed logged the same losses to four
decimals and left adapters of different digests), so its exam counts are a new sample of the same measurement.
\end{sloppypar}

\subsection{The test of this procedure}\label{app:repro-vim-test}

\begin{sloppypar}
The procedure was run on 2026-11-27 in an empty directory on one node of the cluster that ran the experiment (GH200, Ubuntu
26.04 aarch64 in a container; rustc 1.100.0, Python 3.14.4, git-lfs 3.7.1, \texttt{hf} 2.4.0), the blocks extracted from
this file's source and run in one shell, unchanged: \texttt{[setup]}, \texttt{[serve-build]}, \texttt{[smoke]} and
\texttt{[analysis]} (the latter over the smoke's runs, \texttt{RUNS} set to them). The \texttt{[serve]} block was not run,
because every card of the job served the live experiment: \texttt{BRAIN\_CHAT\_URL} was left at its default, that node's
router over the running fleet (brain \texttt{b726c045b}, serving Qwen3-4B), and \texttt{BRAIN\_TRAIN\_URL},
\texttt{BRAIN\_CA\_FILE}, \texttt{BRAIN\_KEY\_FILE}, \texttt{BRAIN\_ARTIFACTS} and \texttt{BRAIN\_ADAPTERS} named one of its
card servers and its shared stores; \texttt{SPLINTER\_GIT} and \texttt{SVEN\_GIT} were the repositories' SSH addresses. Its
two \texttt{brain pull} lines were run alone: every weight file they fetched is byte-identical to the run's model store (3
files of Qwen3-4B, 66 of Qwen3.8-27B-FP8). The \texttt{[full]} blocks were not run; every recipe they use passed
\texttt{splinter validate}.
\end{sloppypar}

\begin{sloppypar}
Wall clock, with a warm Cargo registry and 48 build jobs: clones and builds 8~min (sven 2.1, splinter 2.1, the worker 0.4,
brain 3.4); data 17~s; the task repository and its check 25~s; smoke: forward material and its four datasets 6.0~min
(the task image built from the Debian snapshot included), outcome-first 1.4~min, arm~B 20.6~min, arm~B on its published
dataset 53~min (its training waited behind the shared card's own jobs); analysis and the paper under a minute; 1~h~30~min in
all, 16~GB on disk. The paper block failed in that run, because main's paper had moved to a later evidence revision; it
was changed to read the revision from the paper's own \texttt{data/artifacts.csv} and its commands were then run in the same
directory, building the PDF with the paper's data unchanged. The test found one defect, fixed before this run in splinter \texttt{bde8e09f} (its knowledge entry 233):
an earlier test run, stopped by a time limit while evaluating, could not be continued, because the stopped worker's sandboxes
stayed in the batch directory the controller publishes. With the fix, a smoke run stopped during its evaluations, first by
SIGTERM and then by SIGKILL of its whole process group, was continued by running the smoke block again and committed every
stage. Two trainings of the same records on the same server with the same seed logged the same losses to four decimals and
left adapters of different digests.
\end{sloppypar}
